\section{Methodology}\label{sec:method}

\subsection{What is learned (RQ1)}\label{sec:method-traj}
RQ1 describes each model's representation on three levels. At the level of
the whole space we ask how many dimensions are used and whether the space has
collapsed. At the level of the case we ask what shape the trajectories take.
At the level of progress we ask how step size and turning change as a case
advances. We then compare these profiles across architectures, objectives and
logs.

With $\Delta_t=\z_t-\z_{t-1}$, each case gets a path length
$L=\sum_t\|\Delta_t\|$ and a straightness $S=\|\z_{T'}-\z_1\|/L\in[0,1]$. Each
step gets a step size $v_t=\|\Delta_t\|$, a turning angle
$\theta_t=\arccos\big(\Delta_t^{\top}\Delta_{t+1}/\|\Delta_t\|\|\Delta_{t+1}\|\big)$
and an acceleration $a_t=\|\Delta_{t+1}-\Delta_t\|$. Degenerate inputs return
\texttt{NaN}, never a value that looks good: a collapsed trajectory gives an
undefined $S$, not $S=1$. High $S$ can arise from collapse or from short
trajectories, so trajectory metrics are interpreted only alongside
diagnostics computed on all test representations of a model:
\begin{itemize}[leftmargin=1.2em,itemsep=0pt,topsep=1pt]
  \item total variance;
  \item effective rank $\erank=\exp(-\sum_i p_i\log p_i)$, with $p_i=\lambda_i/\sum_j\lambda_j$;
  \item participation ratio $\PR=(\sum_i\lambda_i)^2/\sum_i\lambda_i^2$;
  \item the distribution of pairwise distances.
\end{itemize}
A model is flagged as collapsed if its variance is at most $10^{-9}$ or
$\erank<1.5$. Because $d$ differs across models, we compare scale-free
quantities only: $\erank/d$, $\PR$, $S$, $\theta_t$, and step sizes normalised
by each model's median pairwise distance. Progress profiles average $v_t$ and
$\theta_t$ within ten bins of normalised progress $s=t/T$. Next-event
trajectories are one prefix shorter than suffix trajectories, so cross-objective
comparisons drop the final suffix prefix.
\QUESTION{Should RQ1 also test what is \emph{decodable} from $\z_t$? Linear
probes (logistic or ridge regression on frozen $\z_t$, fitted on training
prefixes and scored on test) could predict last activity, prefix length,
remaining time and next activity. They would compare against the same probes
on an untrained encoder and on non-learned history encodings. This would answer
``what do the models learn'' directly rather than only through geometry. It
costs about 0.3 pages and needs training-split embeddings, which Phase 7 needs
anyway. Not in the plan yet: parked in PLAN.md (``Parked options'', P1).}

\subsection{Reliability (RQ2)}\label{sec:method-rel}
For every test prediction we fit three nested logistic models of the
probability of an error:
\begin{itemize}[leftmargin=1.2em,itemsep=0pt,topsep=1pt]
  \item[(A)] confidence only: maximum probability, entropy and margin;
  \item[(B)] (A) plus prefix length;
  \item[(C)] (B) plus geometric features of the trajectory so far.
\end{itemize}
The geometric features in (C) are straightness to date, the most recent and
the maximum turning angle, the most recent step size, cumulative path length,
mean acceleration, and the distance to the nearest training representation.
Models are fitted on the validation split and evaluated on test with AUROC,
AUPRC, Brier score and log-likelihood gain. Paired bootstrap confidence
intervals compare C with B. We report high-confidence errors separately,
since that is where confidence alone fails. H1 is supported only if C
improves on B with an interval that excludes zero.
\TODO{Phase 7 not yet run. Validation- and training-split embeddings must first be
extracted, since only test embeddings exist (STATUS next step 9). Fix the
confidence threshold for the ``high-confidence'' subset and the number of
bootstrap resamples.}

\subsection{Future organisation (RQ3)}\label{sec:method-future}
For a fixed-seed sample of $n=500$ test prefixes we build:
\begin{itemize}[leftmargin=1.2em,itemsep=0pt,topsep=1pt]
  \item the latent distance matrix $D_Z$ (Euclidean);
  \item the history distance $D_H$, the normalised Damerau--Levenshtein (edit)
        distance between observed prefixes;
  \item future distances $D_F$ between observed suffixes, recomputed from the
        raw log so that next-event and suffix models are judged against the
        same reference: edit distance, Jaccard distance on activity sets and on
        2-gram sets, and normalised remaining-time difference.
\end{itemize}
We report global alignment $\rho(D_Z,D_F)$ (Spearman), precision@10 of
nearest-neighbour sets, and trustworthiness and
continuity~\cite{venna2001neighborhood}, computed directly on the two distance
matrices. Histories and futures of a case are correlated, so global alignment
alone cannot support H3. H3 rests on two \emph{history-controlled} tests,
with $\eta$ the median of $D_H$:
\begin{enumerate*}[label=(\roman*)]
\item $\rho(D_Z,D_F)$ restricted to pairs with $D_H>\eta$, where a positive
      value means similar futures are retrieved despite dissimilar histories;
\item the same restricted to pairs with $D_H\le\eta$, where a positive value
      means similar histories with divergent futures are pulled apart.
\end{enumerate*}
For H2, each prefix is assigned normalised progress $s\in[0,1]$, binned into
ten bins. Within each bin we compute a Fisher ratio of between-group to
within-group scatter, grouping prefixes by the terminal activity of the case.
A separation ``horizon'' is reported only if it beats a label-permutation
null. Non-learned history encodings (activity counts, last activity) serve
as a floor for every RQ3 statistic.

\subsection{Validation and experiment families}\label{sec:method-families}
Before any real embedding was analysed, every metric was tested (44 unit
tests) on synthetic trajectories with analytically known values:
\begin{itemize}[leftmargin=1.2em,itemsep=0pt,topsep=1pt]
  \item a straight line, with $S=1$ and $\theta=0$;
  \item a circular arc, checked against closed-form values;
  \item a backtracking path, with a single turn of $\theta=\pi$;
  \item a collapsed trajectory, which must return \texttt{NaN};
  \item branching and reconverging ensembles.
\end{itemize}
The history-controlled tests are also checked against a hand-built scenario
with known correlations. Comparisons are split into two families that are
never pooled. \famA (ecological) contains published models, reproduced under
a shared scope; its comparisons are descriptive. \famB (controlled) contains
one encoder trained under the next-event objective (B1) and the suffix
objective (B2), with identical architecture, dimension, optimiser and budget.
B2 adds only a decoder. A difference seen in \famA between next-event and
suffix models is attributed to the objective only if \famB shows it too.
